\documentclass[11pt]{article}
\usepackage[margin=1in]{geometry}
\usepackage[T1]{fontenc}
\usepackage[utf8]{inputenc}
\usepackage{amsmath,amssymb,amsthm}
\usepackage[hidelinks]{hyperref}
\newtheorem{theorem}{Theorem}
\newtheorem{lemma}[theorem]{Lemma}
\newtheorem{corollary}[theorem]{Corollary}
\newcommand{\CPsh}{\operatorname{CPsh}}
\newcommand{\DD}{d'd''}
\title{Curve restrictions of generated weights\\on analytic torus quotients}
\author{Working paper; independent review pending}
\date{9 October 2026}
\begin{document}
\maketitle
\begin{abstract}
We test the entire continuous plurisubharmonic sheaf, including its later
uniform-limit stages, by algebraic analytic curves after algebraically
closed complete field extensions. The proof uses smooth pullback and the
precise curve theorem of Wanner, followed by the minimal-sheaf property.
Decreasing bounded presentations give harmonic comparison on every
strictly affinoid curve domain even when the limit has poles. This supplies
one recognition direction for generated finite-energy weights on the
concrete nonalgebraizable torus quotient. It does not establish the converse
recognition, compactness, envelope regularity, or the full equation theorem.
\end{abstract}

\section{Spaces, weights, and the precise curve tests}
Let $K$ be a complete nontrivially valued non-Archimedean field. Let $X$
be a good separated $K$-analytic space covered by ordinary open sets
isomorphic to ordinary opens in algebraic tori. The running example is
\[
 K=\mathbb C((s)),\qquad |s|=e^{-1/8},\qquad
 X=((\mathbb G_m)^2)^{\mathrm{an}}/
       \langle(s^8,2),(3,s^8)\rangle .
\]
Its previously stated quotient and nonalgebraization foundations remain
separate geometric inputs. This paper uses its open torus charts.
Write $x_i=-\log|z_i|$. On a quotient chart choose the smooth potential
$\rho=\frac12(x_1^2+x_2^2)$ for
$\theta=\DD\rho$. Affine changes of $\rho$ on chart overlaps do not
affect this form.

On any open chart $V$, $\CPsh_V$ is the smallest sheaf of finite continuous
real functions containing ordinary smooth functions with positive
$(1,1)$-Hessian and closed under uniform limits. We retain middle positivity,
which for $(1,1)$-forms is the positive semidefinite cone.
Define $\mathcal C_\theta$ by requiring $\rho+u\in\CPsh_V$ locally.
A generated weight is a pointwise decreasing limit
\[
 u=\lim_j u_j,\qquad u_j\in\mathcal C_\theta.
\]
The generated finite-energy class further requires a common lower bound
for the energies of this presentation. The curve assertion below already
holds without that additional energy bound.

A test consists of a complete algebraically closed nontrivially valued
extension $L/K$, a smooth proper algebraic curve $C/L$, an ordinary open
$W\subset C^{\mathrm{an}}$, and an $L$-analytic map
$f:W\to X_L$. Tests are local on $W$ and on the torus charts of $X_L$.
Here $X_L$ denotes analytic base change and a function on $X$ is pulled
back by the projection $X_L\to X$.
No assertion about every possible nonalgebraic analytic test curve is
included in this definition.

\section{The applied curve theorem and smooth pullback}
The applied published theorem is Wanner's Theorem4.7 in
arXiv:1801.04713v3. With her standing hypotheses, $L$ is algebraically
closed, complete, nontrivially valued, $C$ is smooth proper algebraic,
$W\subset C^{\mathrm{an}}$ is open, and $h:W\to\mathbb R$ is continuous.
If $\DD[h]$ is a positive current, then $h$ is subharmonic in Thuillier's
sense. Her Definition2.4 means that $h$ is upper semicontinuous, is not
identically $-\infty$ on a connected component, and satisfies
\[
 h|_{\partial A}\le H|_{\partial A}\quad\Longrightarrow\quad h|_A\le H
 \tag{1}
\]
for every strictly $L$-affinoid curve domain $A\subset W$ and every
harmonic function $H$ on $A$. Such harmonic functions are continuous;
adding a real constant preserves harmonicity.
Finite continuous functions automatically satisfy the condition excluding
an identically infinite value.

We use the ordinary smooth-form pullback rules
\[
 \DD(h\circ g)=g^*(\DD h),\qquad
 g^*(\hbox{positive $(1,1)$-form})\ \hbox{is positive}. \tag{2}
\]
They concern smooth forms and actual analytic morphisms. Locally they are
the smooth chain rule in unit-moment presentations: compose the analytic
units with the morphism, refine the moment list, and pull back the real
positive semidefinite Hessian. A rank-one positive covector form pulls
back to the corresponding positive rank-one form; sums and limits give
the positive semidefinite assertion. Field extension uses the same unit
expressions and real coefficient functions. A positive smooth form gives
a positive current under canonical curve integration. These are the
ordinary CLD form and positivity rules, also used in Wanner's
Proposition3.20 and Remark3.21. We take no pullback of an arbitrary
positive current or of an arbitrary distributional-psh function.

Consequently every ordinary smooth psh function on a torus chart pulls
back under a test to a finite smooth function on a curve with positive
Hessian current. Wanner's theorem applies locally. The target chart is an
open in an algebraic torus, so the algebraic-chart scope of her smooth
pullback remark suffices; the global quotient need not be algebraizable.

\section{Uniform limits and the entire continuous sheaf}
\begin{lemma}\label{uniform}
Suppose finite continuous subharmonic functions $h_j$ on $W$ converge
locally uniformly to a finite continuous $h$. Then $h$ satisfies (1).
\end{lemma}
\begin{proof}
Fix $A,H$ as in (1) and suppose $h\le H$ on $\partial A$.
Compactness of $A$ gives uniform convergence there. For every $\epsilon>0$
and every sufficiently large $j$, $h_j\le H+\epsilon$ on $\partial A$.
Comparison for $h_j$ gives $h_j\le H+\epsilon$ on $A$. Take the limit and
then let $\epsilon$ decrease to zero. The remaining defining conditions
hold because $h$ is finite and continuous.
\end{proof}

\begin{theorem}\label{recognition}
For every $h\in\CPsh_V$ and every test mapping into $V_L$, the pullback
$h\circ f$ is a finite continuous subharmonic function on $W$.
This applies to every stage of the minimal continuous-psh sheaf.
\end{theorem}
\begin{proof}
On the opens of $V$, let $\mathcal Q$ consist of finite continuous
functions whose pullback under every specified curve test is subharmonic.
Locality of the curve definition, or equivalently its sheaf property
in Wanner's Proposition2.5, makes $\mathcal Q$ a sheaf.
The smooth pullback calculation and Theorem4.7 show that it contains the
ordinary smooth psh functions. It is closed under uniform limits:
composition preserves a uniform bound, and Lemma\ref{uniform} applies
on each compact test domain. Hence $\mathcal Q$ has all the closure
properties in the definition of $\CPsh_V$. Minimality gives
$\CPsh_V\subset\mathcal Q$.
\end{proof}
This argument does not identify later sheaf stages with locally uniformly
smooth-approximable functions. It is a recognition by a new closed sheaf.

\section{Decreasing limits, poles, and affine-boundary comparison}
\begin{lemma}\label{decreasing}
Let finite continuous subharmonic functions $h_j$ decrease on $W$ to
$h:W\to[-\infty,\infty)$. Then $h$ is upper semicontinuous and satisfies
(1). On any connected component on which it is not identically
$-\infty$, it is subharmonic in Definition2.4.
\end{lemma}
\begin{proof}
The decreasing infimum of the continuous functions is upper
semicontinuous. Fix $A,H$ and assume $h\le H$ on $\partial A$.
For $\epsilon>0$ the relatively open sets
\[
 U_j=\{y\in\partial A:h_j(y)<H(y)+\epsilon\}
\]
increase and cover the compact set $\partial A$.
They cover points at which $h=-\infty$ as well as finite points.
A finite subcover and monotonicity give a single index $j_0$ such that
$h_j<H+\epsilon$ everywhere on $\partial A$ for all $j\ge j_0$.
Thus (1) for $h_j$ gives $h_j\le H+\epsilon$ on $A$.
Pass to the decreasing limit and let $\epsilon$ decrease to zero.
The assertion on a component follows directly from the exact definition,
which is the only place its condition excluding an identically
$-\infty$ component is needed.
\end{proof}

\begin{corollary}\label{weights}
Let $u_j\in\mathcal C_\theta$ decrease to $u$. On every test lying over
a chart with reference potential $\rho$, the function
\[
 q=\rho\circ f+u\circ f
\]
is upper semicontinuous and satisfies harmonic comparison (1).
It is subharmonic on each component where it is not identically
$-\infty$. This holds in particular for every generated finite-energy
weight, every decreasing presentation of it, and every test field $L$
specified above.
\end{corollary}
\begin{proof}
The functions $q_j=(\rho+u_j)\circ f$ are finite continuous
subharmonic by Theorem\ref{recognition}. They decrease to $q$.
Apply Lemma\ref{decreasing} on the test region lying in this chart.
\end{proof}

\begin{corollary}\label{maximum}
For a strictly affinoid domain $A$ contained in such a chart test,
\[
 \sup_A q\le\sup_{\partial A}q.
\]
If $q=-\infty$ on the whole boundary, then $q=-\infty$ on $A$.
\end{corollary}
\begin{proof}
Every real constant majorizing the boundary values is harmonic.
Apply (1) and take the infimum of these constants. When all boundary
values are infinite, every real constant qualifies.
\end{proof}

\section{Scope and remaining analytic gates}
The result supplies comparison on the stated algebraic curve tests for
the full continuous-psh sheaf and its decreasing generated weights.
It is one direction of a possible intrinsic recognition theorem.
It proves neither that all curve-test weights admit global bounded
decreasing presentations nor that bounded discontinuous intrinsic
weights belong to the generated class.

In particular curve-wise comparison does not by itself give compactness
of a normalized family of weights on the analytic surface, a global
Shilov-boundary maximum principle in higher dimension, or domination
and uniqueness for its two-dimensional Monge--Ampere operator.
Actual separated compactness, continuous-obstacle envelopes and
orthogonality, arbitrary-support finite-energy measure existence, and
uniqueness modulo constants remain to be proved. Original P16.7 remains
open; this working paper is independently unreviewed.

\section{Source version and proof-reading boundary}
The exact applied Wanner snapshot is arXiv:1801.04713v3, dated
8 November2018, with the algebraically closed nontrivially valued field
convention and proper algebraic curve scope retained above.
Definition2.4/Proposition2.5, Proposition3.20/Remark3.21, and the entire
Lemma4.6/Theorem4.7 statement and immediate proof were read from the
retained full text. The proof includes the compactly supported model test,
its nonnegative smooth approximations and the final positive-current
contradiction. Its cited semistable reduction, model-function and curve
potential-theory backgrounds remain named published foundations; their
entire recursive proofs are not claimed freshly audited here.
The local pullback calculation (2) uses ordinary smooth unit moments
and positive $(1,1)$-forms, with no arbitrary-current pullback theorem.

\begin{thebibliography}{9}
\bibitem{Wanner} V. Wanner,
\emph{Comparison of two notions of subharmonicity on non-Archimedean
curves}, arXiv:1801.04713v3,8 November2018,
\url{https://arxiv.org/abs/1801.04713v3}.
Published in Math. Z.293 (2019),443--474,
DOI10.1007/s00209-018-2205-z. The cited snapshot is not a byte
identification with the published article.
\bibitem{CLD} A. Chambert-Loir and A. Ducros,
\emph{Formes differentielles reelles et courants sur les espaces de
Berkovich}, arXiv:1204.6277v2,25 July2025.
The applied continuous-sheaf convention is Definition8.1.7.
\end{thebibliography}
\end{document}
